\documentclass[../../main.tex]{subfiles}
\begin{document}

\begin{flushright}
\footnotesize\textit{【自動文字起こし・要確認】}
\end{flushright}

{\bf [解]}
$0<p<1$，$q>0\quad\cdots①$

\paragraph{(1)}

①及び$0<x<1$から，$0<e^{-qx}<1$だから，
\begin{align}
0<(1-p)x<x,\qquad0<(1-x)(1-e^{-qx})<1-x
\end{align}
辺々加えて
\begin{align}
0<f(x)<1
\end{align}

\paragraph{(2)}

(1)及び$0<x_0<1$から，帰納的に任意の$n\ge0$に対し$0<x_n<1\quad\cdots②$である．
\begin{align}
p>q\quad\cdots③,\qquad e^x\ge1+x\quad\cdots④
\end{align}
に注意して，$A=e^{-qx_n}$とすると，
\begin{align}
x_{n+1}&=(A-p)x_n+1-A\\
&\le(A-p)x_n+qx_n\quad(\because④より e^{-qx_n}\ge1-qx_n\text{，すなわち}A\ge1-qx_n\text{，つまり}1-A\le qx_n)\\
&=(A-p+q)x_n\\
&<(1-p+q)x_n\quad(\because①より1>A，②より0<x_n)
\end{align}
両辺絶対値をとってくり返し用いると，
\begin{align}
|x_n|<|1-p+q|^nx_0\quad\cdots⑤
\end{align}
③及び①から$0<q<p<1$，よって$0<1-p+q<1$，つまり$|1-p+q|<1$だから，$n\to\infty$の時，⑤の右辺は$0$に収束するので，はさみうちから
\begin{align}
|x_n|\to0\quad\therefore\ x_n\to0
\end{align}

\paragraph{(3)}

$g(x)=f(x)-x$とおく．$0<x<1$かつ$g(x)=0$なる$x$があることを示せば良い．$A=e^{-qx}$として，
\begin{align}
g(x)&=-px+(1-x)(1-A)\\
&=(A-1-p)x+(1-A)\\
&\ge(-qx-p)x+(1-A)\quad(\because④\text{より}A-1\ge-qx)\quad\cdots⑥
\end{align}
である．⑥の右辺を$h(x)$とおく．
\begin{align}
h(x)=-qx^2-px+1-e^{-qx}
\end{align}
\begin{align}
h'(x)&=-2qx-p+qA\\
h''(x)&=-2q-q^2A<0\quad(\because q>0,\ A>0)
\end{align}
から，$h'(x)$は単調減少で，
\begin{align}
h'(0)=q-p>0\quad(\because p<q),\qquad h'(1)=q(e^{-q}-2)-p<0\quad(\because e^{-q}<1<2,\ p>0)
\end{align}
だから，$h'(\alpha)=0$なる$\alpha$が$0<\alpha<1$に存在し，下表を得る（$\because h'$は連続）．
\begin{center}
\begin{tabular}{c|ccccc}
$x$ & $0$ & $\cdots$ & $\alpha$ & $\cdots$ & $1$ \\
\hline
$h'$ & & $+$ & $0$ & $-$ & \\
$h$ & $0$ & $\nearrow$ & & $\searrow$ & \\
\end{tabular}
\end{center}
ここで$h(0)=0$だから，$0<x<\alpha$では$h(x)>0\quad\cdots⑦$である．したがって，⑥⑦から，
\begin{align}
g(x)\ge h(x)>0
\end{align}
なる$x$が$0<x<1$に存在する．$\cdots⑧$

一方，$g(1)=f(1)-1=(1-p)-1=-p<0\quad\cdots⑨$である．⑧⑨及び$g$が連続であることから，$g(x)=0$となる$x$が$0<x<1$に存在する．$\blacksquare$

\end{document}
